% ============================================================
%  Setmana 19 — Unitat 1 · Matrius
%  Exercicis: 42, 49, 69, 81, 85
% ============================================================

\setmana{19}{Unitat 1 \textperiodcentered{} Matrius}{24 de gener de 2027}%
        {42, 49, 69, 81 i 85}

\begin{errada}{Exercici 49 d)}
  \ondiu{$BC=\begin{pmatrix}-7 & 3 & 3\\ -1 & -1 & 8\end{pmatrix}$}
  \hadedir{$BC=\begin{pmatrix}-8 & 3 & 3\\ -4 & -1 & 8\end{pmatrix}$}
  \motiu{Els dos elements de la primera columna estan malament:
    $(-2)\cdot 4+1\cdot 0+1\cdot 0=-8$ i $(-1)\cdot 4+0\cdot 0+3\cdot 0=-4$.}
\end{errada}

%\senseerrades{42, 69, 81 i 85}

% ---------- Observacions (només surten amb \obstrue) ----------
\ifobs
\bigskip
\noindent
Dues observacions, que no són cap errada:\par\nopagebreak

\begin{observacio}{Exercici 42}
  L'enunciat demana dos exemples de cada tipus de matriu i la solució només
  en dona un.
\end{observacio}

\begin{observacio}{Exercici 85 b)}
  La solució no exclou $a=0$ ni $b=0$. Si $a=b=0$, $B$ és la matriu nul·la, que no és nilpotent de grau 2 perquè ja és nul·la amb l'exponent 1. Cal afegir-hi $a\neq 0$ a la primera família de matrius i $b\neq 0$ a la segona.
\end{observacio}
\fi
